\section{Exact ranks and Smith forms for the level-four product system}
\label{newrank:sec:main}

The manuscript's odd-weight rank conjecture concerns a fully specified
matrix of depth-two relations.  It does not concern the dimension of the
space of numerical periods.  We prove that conjecture, determine the
corresponding matrices in even weight, and compute their integral Smith
forms.  The reason for the simple answer is a polynomial involution which
becomes a permutation of monomials after an integral change of variables.

\subsection{The matrix and the result}

For $a+b=w$, write formal coordinates
\[
 x_{a,b}^{r,s},\qquad r,s\in\mathbb Z/4\mathbb Z,
\]
representing $\Im\operatorname{Li}_{a,b}(i^r,i^s)$.  Set them equal to zero
when $r,s$ are both even, and impose
$x_{a,b}^{r,s}=-x_{a,b}^{-r,-s}$.  The six canonical color pairs are
\begin{equation}\label{newrank:eq:pairs}
 (0,1),\ (1,0),\ (1,1),\ (1,2),\ (1,3),\ (2,1).
\end{equation}
Remove $x_{1,w-1}^{0,1}$, the single divergent imaginary coordinate after
conjugation.  Thus the coordinate count is $6w-7$.

For $p+q=w$ and $x=i^r,y=i^s$, the two product rows have left sides
\begin{align}
 \mathsf S_{p,q}^{r,s}
 &=x_{p,q}^{r,s}+x_{q,p}^{s,r},\label{newrank:eq:stuffle}\\
 \mathsf H_{p,q}^{r,s}
 &=\sum_{j=0}^{p-1}\binom{q+j-1}{j}x_{q+j,p-j}^{s,r-s}
  +\sum_{j=0}^{q-1}\binom{p+j-1}{j}x_{p+j,q-j}^{r,s-r}.
 \label{newrank:eq:shuffle}
\end{align}
Their analytic right sides are respectively
\[
 \Im\bigl(\operatorname{Li}_p(x)\operatorname{Li}_q(y)
           -\operatorname{Li}_w(xy)\bigr),\qquad
 \Im\bigl(\operatorname{Li}_p(x)\operatorname{Li}_q(y)\bigr).
\]
When both single factors are finite, retain both rows.  When exactly one
factor is $\operatorname{Li}_1(1)$, retain the difference
$\mathsf S-\mathsf H$, in which the divergent coordinate cancels.
Its right side is $-\Im\operatorname{Li}_w(xy)$.
Discard zero rows and retain the remaining integer coefficient matrix
$A_w$.  Repeated rows do not affect the assertions below.  This is
precisely the matrix in the manuscript's Conjecture~\texttt{signed:conj:rank}.
Let $A_w^{\neg g}$ be the matrix obtained by deleting the $w-1$ columns
$x_{a,w-a}^{1,0}$; it has $5w-6$ columns.

\begin{theorem}[All-weight ranks and integral invariant factors]
\label{newrank:thm:smith}
For every odd $w\ge3$,
\begin{equation}\label{newrank:eq:odd}
 \operatorname{rank}A_w=5w-6,\qquad
 \operatorname{rank}A_w^{\neg g}=\frac{9w-11}{2}.
\end{equation}
Every nonzero Smith invariant of either integer matrix is one.
For every even $w\ge4$, both matrices have full column rank, and their
nonzero Smith invariants are
\begin{align}
 \operatorname{SNF}(A_w)
 &=\operatorname{diag}\bigl(1^{[5w-7]},2^{[w]}\bigr),
 \label{newrank:eq:evenfull}\\
 \operatorname{SNF}(A_w^{\neg g})
 &=\operatorname{diag}\bigl(1^{[(9w-12)/2]},2^{[w/2]}\bigr).
 \label{newrank:eq:evennong}
\end{align}
Here the bracketed exponent denotes multiplicity; extra zero rows in a
rectangular Smith form are suppressed.
\end{theorem}

\subsection{A unit pivot restores the omitted coordinate}

Restore $t=x_{1,w-1}^{0,1}$ and impose both homogeneous rows
\eqref{newrank:eq:stuffle}--\eqref{newrank:eq:shuffle} for every color and
index pair.  Denote the resulting auxiliary matrix by $M_w$.
The excluded stuffle row for $(p,q,r,s)=(1,w-1,0,1)$ is
\begin{equation}\label{newrank:eq:restore}
 t+x_{w-1,1}^{1,0}=0.
\end{equation}
Every excluded stuffle row with a nonzero imaginary part is this row up
to sign and interchange of factors.  The other excluded stuffle rows
vanish under the imposed color conventions.  Adding
\eqref{newrank:eq:restore} to the baseline difference rows therefore
recovers all the missing individual stuffle and shuffle rows by integral
row combinations.  Conversely, their differences recover the baseline
rows.  A finite-product row cannot contain $t$: an outer index one and
outer color zero in a shuffle term require the corresponding input
factor to be $\operatorname{Li}_1(1)$.

The coefficient of $t$ in \eqref{newrank:eq:restore} is one, so eliminating
it is unimodular.  The row-module quotients for $M_w$ and $A_w$ are
isomorphic; the Smith form of $M_w$ has exactly one additional unit.
Their homogeneous nullspaces over any field are isomorphic, with
$t=-x_{w-1,1}^{1,0}$.  After deleting the Gaussian columns, the same
argument applies with those coordinates set to zero and $t=0$.

This auxiliary construction is a statement about integer coefficient
matrices.  It assigns no analytic finite value to the divergent
coordinate and introduces no new analytic relation.

\subsection{Six polynomials reduce to two involutions}

Put $n=w-2$ and encode a vector of coordinates by homogeneous polynomials
\[
 f_{r,s}(X,Y)=\sum_{a=1}^{w-1}
 x_{a,w-a}^{r,s}X^{a-1}Y^{w-a-1}.
\]
The complete homogeneous product equations are exactly
\begin{align}
 f_{r,s}(X,Y)+f_{s,r}(Y,X)&=0,\label{newrank:eq:polyS}\\
 f_{s,r-s}(X+Y,X)+f_{r,s-r}(X+Y,Y)&=0.
 \label{newrank:eq:polyH}
\end{align}
For example, the coefficient of $X^{p-1}Y^{q-1}$ in the second term of
\eqref{newrank:eq:polyH} is
\[
 \sum_{j=0}^{q-1}\binom{p+j-1}{j}x_{p+j,q-j}^{r,s-r},
\]
which is the second binomial sum in \eqref{newrank:eq:shuffle}.

Write
\[
 A=f_{0,1},\quad B=f_{1,0},\quad C=f_{1,1},\quad
 D=f_{1,2},\quad E=f_{1,3},\quad F=f_{2,1}.
\]
The stuffle equations become
\begin{align}
 A(X,Y)&=-B(Y,X),&F(X,Y)&=-D(Y,X),\label{newrank:eq:AF}\\
 C(X,Y)&=-C(Y,X),&E(X,Y)&=E(Y,X).
 \label{newrank:eq:CE}
\end{align}
Up to conjugation and interchange of the factors, the nonzero shuffle
equations are exhausted by colors $(0,1),(1,1),(2,1),(1,3)$.
With $V=X+Y$, their left sides are respectively
\[
 E(V,X)+A(V,Y),\quad B(V,X)+B(V,Y),\quad
 C(V,X)-F(V,Y),\quad -D(V,X)+D(V,Y).
\]
The color pairs $(0,0)$, $(0,2)$, $(2,0)$, $(2,2)$ vanish. Conjugation supplies
the pairs $(0,3)$, $(3,0)$, $(3,3)$, $(3,2)$, $(3,1)$, $(2,3)$; the remaining pairs
$(1,0)$ and $(1,2)$ are factor interchanges.  This checks all sixteen
color pairs.

Combining these equations with \eqref{newrank:eq:AF} gives
\begin{align}
 E(X,Y)&=B(X-Y,X),& B(X,X-Y)&=-B(X,Y),\label{newrank:eq:B}\\
 C(X,Y)&=-D(X-Y,X),& D(X,X-Y)&=D(X,Y).
 \label{newrank:eq:D}
\end{align}
Thus $A,F,E,C$ are eliminated and no relation couples $B$ to $D$.

Let $T(X,Y)=(X,X-Y)$.  The remaining conditions
\eqref{newrank:eq:CE}, transported through the substitution
$U(X,Y)=(X-Y,X)$, yield the coefficient rows
\begin{equation}\label{newrank:eq:normal}
 (I+T)B,\quad (I-(-1)^nT)B,\qquad
 (I-T)D,\quad (I+(-1)^nT)D.
\end{equation}
Here $T$ acts by substitution.  To check the signs, if $S(X,Y)=(Y,X)$,
then $USU^{-1}=-T$, and a degree-$n$ homogeneous polynomial takes the
factor $(-1)^n$ under simultaneous sign change.

These are integral presentation operations.  Each of the four eliminated
polynomials has coefficient one.  The substitutions and their inverses
have integer coefficients and determinant $\pm1$, so they induce
unimodular maps of the coefficient lattices.  Consequently the four
eliminations contribute $4(w-1)$ Smith units, and the remaining Smith
problem is exactly \eqref{newrank:eq:normal}.

\subsection{The integral monomial permutation proves the theorem}

Make the unimodular change of variables
\begin{equation}\label{newrank:eq:uv}
 u=Y,\qquad v=X-Y,\qquad X=u+v,\quad Y=u.
\end{equation}
The involution $T$ now exchanges $u,v$.  It therefore permutes the
integral monomial basis $u^jv^{n-j}$ by $j\leftrightarrow n-j$.

If $n=2m-1$ is odd, this permutation has $m$ two-element orbits and no
fixed monomial.  On each orbit the two $B$ relations in
\eqref{newrank:eq:normal} coincide and give the primitive row $(1,1)$.
The two $D$ relations coincide and give $(1,-1)$.  Each orbit contributes
one Smith unit and one free coordinate in each polynomial.  Hence the
full nullity is $2m=w-1$, and the nullity with $B=0$ is $m=(w-1)/2$.
Counting the $4(w-1)$ eliminated units and subtracting the auxiliary
unit \eqref{newrank:eq:restore} gives respectively
\[
 4(w-1)-1+(w-1)=5w-6,
 \qquad
 4(w-1)-1+\frac{w-1}{2}=\frac{9w-11}{2}
\]
nonzero Smith units.  This proves \eqref{newrank:eq:odd} and its
integral strengthening.

If $n=2h$ is even, there are $h$ two-element orbits and one fixed
monomial.  Both plus and minus relations occur for each of $B,D$.
A two-element orbit gives
\[
 \begin{pmatrix}1&1\\1&-1\end{pmatrix},
\]
with Smith form $\operatorname{diag}(1,2)$; a fixed monomial gives the
single nonzero row $(2)$.  Each polynomial therefore contributes $h$
units and $h+1$ twos, with no free coordinate.  Adding the eliminated
units and removing the auxiliary unit gives
\[
 4(w-1)-1+(w-2)=5w-7\quad\hbox{units},\qquad w\quad\hbox{twos}
\]
for $A_w$, and
\[
 4(w-1)-1+\frac{w-2}{2}=\frac{9w-12}{2}\quad\hbox{units},
 \qquad\frac w2\quad\hbox{twos}
\]
for $A_w^{\neg g}$.  This proves
\eqref{newrank:eq:evenfull}--\eqref{newrank:eq:evennong} and completes
the proof of Theorem~\ref{newrank:thm:smith}.

\subsection{Primitive witnesses and exact certificate consequences}

For odd $n=2m-1$ and $0\le j<m$, define
\begin{align*}
 P_j(X,Y)&=Y^j(X-Y)^{n-j}-Y^{n-j}(X-Y)^j,\\
 Q_j(X,Y)&=Y^j(X-Y)^{n-j}+Y^{n-j}(X-Y)^j.
\end{align*}
The complete integer nullspace is obtained, uniquely, by setting
\[
 B=\sum_{j=0}^{m-1}b_jP_j,\qquad
 D=\sum_{j=0}^{m-1}d_jQ_j,\qquad b_j,d_j\in\mathbb Z,
\]
using \eqref{newrank:eq:AF}, \eqref{newrank:eq:B},
\eqref{newrank:eq:D} for the other colors, and omitting the restored
coordinate.  Primitivity follows directly from the monomial pairs in
the unimodular coordinates \eqref{newrank:eq:uv}.

\begin{corollary}[All Gaussian-only consequences]
\label{newrank:cor:gaussian}
In odd weight $w$, the rational row space of $A_w$ supported only on the
Gaussian columns has dimension $(w-1)/2$.  It is exactly the span of the
$(w-1)/2$ same-point shuffle rows.  Every integer vector in this space
is a unique integer combination of those rows.
\end{corollary}
\begin{proof}
The kernel of projection of the row space to non-Gaussian columns has
dimension
$\operatorname{rank}A_w-\operatorname{rank}A_w^{\neg g}=(w-1)/2$.
The same-point shuffle rows with $p+q=w$, $1\le p\le m=(w-1)/2$,
have Gaussian coefficients
\[
 R_{p,a}=\binom{a-1}{p-1}+\binom{a-1}{w-p-1}.
\]
Their first $m$ columns form the triangular Pascal matrix
$\binom{a-1}{p-1}$, with determinant one.  They are independent and
span the indicated rational space.  The same unit minor shows that an
integer vector in this space has integer expansion coefficients.
\end{proof}

The Smith forms have two further direct consequences.  In odd weight,
every integer coefficient vector in the rational row span has an integer
row certificate.  In even weight, twice every coordinate vector has an
integer row certificate.  Thus every convergent imaginary level-four
double in even weight reduces with coefficients in $\tfrac12\mathbb Z$
to the raw finite-product and single-value right sides of the stated
system.  This denominator assertion is relative to that raw vocabulary;
rewriting the single values in powers of $\pi$, zeta values, and beta
values can introduce their usual additional rational factors.

At weight five a particularly short obstruction is obtained by choosing
$B=0,D=X^3$.  Then
\[
 C=-(X-Y)^3,\quad F=-Y^3,\quad A=E=0.
\]
The six nonzero coordinates of this integer kernel vector are
\[
 x_{4,1}^{1,2}=1,\quad x_{1,4}^{2,1}=-1,\quad
 x_{4,1}^{1,1}=-1,\quad x_{3,2}^{1,1}=3,\quad
 x_{2,3}^{1,1}=-3,\quad x_{1,4}^{1,1}=1.
\]
The manuscript's target
$3g_{4,1}+3g_{3,2}+9g_{2,3}+7u_{4,1}$ pairs with this vector to give
$7$.  This proves again that the proposed mixed relation is absent from
the specified linear product system.  It does not disprove the numerical
identity or preclude its derivation from additional relations.

\paragraph{Scope and verification.}
The theorem proves the manuscript's conjecture and a stronger integral
normal form.  It makes no assertion that this relation system contains
all standard polylogarithm relations, and no assertion of arithmetic
independence.  Formal double spaces, polynomial relation methods, and
colored regularization are established tools; appropriate context is
provided by Yuan--Zhao \cite{YuanZhao} and Zhao \cite{Zhao}.  The new contribution here
is the exact normal form of the particular matrix family under study.
The independent companion script reconstructs the displayed rows,
recovers the baseline $92\times23$ matrix size at weight five, checks
ranks and primitive nullspace residuals through weight thirteen, and
checks both Smith forms through weight twelve.  The general result rests
on the integral proof above, rather than on extrapolation from these
finite checks.
